\chapter{2023年「线性代数」\titlelogo{https://sci.hdu.edu.cn}{HDU}期中回忆}
\date{2023年12月3日}{基础数学教学团队}{A0714030}
{未央数学讲师~xxy}
{\qrcode{https://qm.qq.com/q/OlkSiy9W04}}{\href{https://qm.qq.com/q/OlkSiy9W04}{HDU 数学营}}
{\qrcode{http://weixin.qq.com/r/mhKYgFTEMQlOrRDj90eI}}{\href{http://weixin.qq.com/r/mhKYgFTEMQlOrRDj90eI}{未央学社公众号}}

\section{填空题（每题3分，共18分）}
\begin{problem}
    设矩阵$\mat{A}=
    \begin{pmatrix}
        1 &-1 & 2\\
        0 & 2 & a\\
        2 & 1 & 5
    \end{pmatrix}$
    为满秩矩阵，则$a$的取值为\ans{$a\neq\frac23$}.
\end{problem}
\begin{solution}*
    $\abs{\mat{A}}=2-3a\neq 0$，即$a\neq\frac23$，此时$\mat{A}$满秩.
\end{solution}

\begin{problem}
    若$\text{R}(\mat{A}_{5\times 4})=4,\ \text{R}(\bf{B}_{4\times 5})=3$，则$\text{R}(\bf{AB})=$\ans{$3$}.
\end{problem}
\begin{solution}*
    可逆矩阵又称满秩矩阵. $\mat{A}_{5\times 4}$列满秩，所以$\rank{\mat{A}_{5\times 4}\mat{B}_{4\times 5}}=\rank{\mat{B}}=3$.
\end{solution}

\begin{problem}
    已知$\mat{A}=
    \begin{pmatrix}
        2 & -1\\
        -3 & 3\\
    \end{pmatrix},f(x)=x^2-5x+4,\text{则}f(\mat{A})$=\ans{
    $\begin{pmatrix}
        1 & 0\\
        0 & 1
    \end{pmatrix}$}.
\end{problem}
\begin{solution}*
    $f(\mat{A})=(\mat{A}-\mat{E})(\mat{A}-4\mat{E})=\begin{pmatrix}
        1 & -1\\-3 &2
    \end{pmatrix}
    \begin{pmatrix}
        -2 & -1\\
        -3 & -1
    \end{pmatrix}=
    \begin{pmatrix}
        1 & 0\\
        0 & 1
    \end{pmatrix}$.
\end{solution}
    
\begin{problem}
    设四阶矩阵$\mat{A}=\left(\alpha,\gamma_{2},\gamma_{3},\gamma_{4}\right),\mat{B}=\left(\beta,\gamma_{2},\gamma_{3},\gamma_{4}\right),\text{其中}\alpha,\beta,\gamma_{2},\gamma_{3},\gamma_{4}$均为4维列向量. 已知$\abs{\mat{A}}=4,\ \abs{\mat{B}}=2$，则$\abs{\mat{A}+\mat{B}}=$\ans{48}.
\end{problem}
\begin{solution}*
    $\mat{A}+\mat{B}=(\alpha+\beta,2\gamma_2,2\gamma_3,2\gamma_4),\ \abs{\mat{A}+\mat{B}}=2^3(\abs{\mat{A}}+\abs{\mat{B}})=8\times(4+2)=48$.
\end{solution}

\begin{problem}
已知四阶矩阵$\mat{A}$的第一行元素依次为$-1,0,2,4$，$\abs{\mat{A}}$中第二行元素的余子式分别为$10,5,a,2$则$a=$\ans{9}.
\end{problem}
\begin{solution}
    第二行代数余子式为$-10,\ 5,\ -a,\ 2$，由代数余子式性质知$(-1)\times(-10)+0\times 5+2\times(-a)+4\times2=0$，即$a=9$.
\end{solution}

\begin{problem}
    设三阶矩阵$\mat{A}$=$\begin{pmatrix}
        a & b & b\\
        b & a & b\\
        b & b & a
    \end{pmatrix}$,若$\mat{A}$的伴随矩阵的秩等于$1$，则必有\ans{$a=-2b$且$a\neq b$}（填$a$与$b$的关系）.
\end{problem}
\begin{solution}
\begin{itemize}
    \item 将$\mat{A}$化为行阶梯形矩阵
    \[\mat{A}=\begin{pmatrix}
        a & b & b\\
        b & a & b\\
        b & b & a
    \end{pmatrix}\overset{c}{\sim}
    \begin{pmatrix}
        a+2b & b & b\\
        a+2b & a & b\\
        a+2b & b & a
    \end{pmatrix}\overset{r}{\sim}
    \begin{pmatrix}
        a+2b & b & b\\
        0 & a-b & 0\\
        0 & 0 & a-b
    \end{pmatrix}\]
    \item $\mat{A}$的秩为$2$，所以$(a+2b)(a-b)^2=0$，得$a=-2b$或$a=b$.
    \item 若$a=b$，则$\rank{\mat{A}}=1$与题目矛盾. 故$a=-2b$且$a\neq b$.
\end{itemize}
\end{solution}

\section{选择题（每题3分，共18分）}
\setcounter{prob}{0}

\begin{choice}{A}
设$\mat{A}=\begin{pmatrix}
    1 & 2\\
    -3 & 4
\end{pmatrix},\ \mat{B}=
\begin{pmatrix}
    \frac25 & x  \\
    \addlinespace[6pt]
    \frac{3}{10} & y 
\end{pmatrix}$，
若$\mat{B}$是$\mat{A}$的逆矩阵，则$x=$
\begin{tasks}(4)
    \task $-\frac15$
    \task $\frac15$
    \task $\frac{1}{10}$
    \task $-\frac{1}{10}$
\end{tasks}
\end{choice}
\begin{solution}
    $\abs{\mat{A}}=10$，所以$\mat{A}^{-1}=\frac{1}{\abs{\mat{A}}}\begin{pmatrix}
        \mathrm{A}_{11} & \mathrm{A}_{21}\\
        \mathrm{A}_{12} & \mathrm{A}_{22}
    \end{pmatrix},\ x=\frac{\mathrm{A}_{21}}{\abs{\mat{A}}}=-\frac15$. \sokka{A}
\end{solution}

\begin{choice}{B}
    设$\mat{A}$为$n$阶矩阵，则$\abs{\mat{A}}=0$的必要条件是
\begin{tasks}(2)
    \task $\abs{\mat{A}}$中有两行（列）元素对应成比例
    \task $\abs{\mat{A}}$中必有一行为其余行的线性组合
    \task $\abs{\mat{A}}$中有一行元素全为零
    \task $\abs{\mat{A}}$中任意一行为其余行的线性组合
\end{tasks}
\end{choice}
\begin{solution}*
    考虑逆否命题，若任意一行都不是其余行线性组合，则$\mat{A}$满，$\abs{\mat{A}}\neq0$. \sokka{B}
\end{solution}

\begin{choice}{B}
    设$\mat{A}$为$n$阶矩阵，$\mat{B}$是$n$阶反对称矩阵，则下列矩阵中为反对称矩阵的是
\begin{tasks}(4)
    \task $\mat{BAB}$
    \task $\mat{ABA}^\T$
    \task $\mat{ABAB}$
    \task $\mat{BABA}$
\end{tasks}
\end{choice}
\begin{solution}*
    $(\mat{ABA}^\T)^\T=\mat{AB}^\T\mat{A}^\T$.\sokka{B}
\end{solution}
        
\begin{choice}{C}
    方程组$\mat{AX} = \mat{B}$中，方程的个数小于未知数的个数，则有
\begin{tasks}(4)
    \task $\mat{AX}=\mat{B}$必有无穷多解
    \task $\mat{AX}=\mat{B}$一定无解
    \task $\mat{AX}=0$必有非零解
    \task $\mat{AX}$
\end{tasks}
\end{choice}
\begin{solution}*
    见线性方程组的解理论第一节.\sokka{C}
\end{solution}

\begin{choice}{C}
    设$\mat{A}=
    \begin{pmatrix}
        2 & 1 & 8 & 3 & 7\\ 
        2 & -3 & 0 & 7 & -5\\
        3 & -2 & 5 & 8 & 0\\
        1 & 0 & 3 & 2 & 0\\
    \end{pmatrix}$，则$\rank{\mat{A}}=$
\begin{tasks}(4)
    \task $1$
    \task $2$
    \task $3$
    \task $4$
\end{tasks}
\end{choice}
\begin{solution}
    \[\mat{A}\overset{r}{\sim}\begin{pmatrix}
        1 & 0 & 3 & 2 & 0\\
        2 & 1 & 8 & 3 & 7\\
        2 & -3 & 0 & 7 & -5\\
        3 & -2 & 5 & 8 & 0
    \end{pmatrix}\overset{r}{\sim}
    \begin{pmatrix}
        1 & 0 & 3 & 2 & 0\\
        0 & 1 & 2 & -1 & 7\\
        0 & -3 & -6 & 3 & -5\\
        0 & 0 & -4 & 2 & 0
    \end{pmatrix}\overset{r}{\sim}
    \begin{pNiceMatrix}
    \CodeBefore
        \tikz \draw [H6,thick] (2-|1) |- (2-|2) |- (3-|2) |- (3-|5) |- (4-|5) |- (4-|6) ;
    \Body
        1 & 0 & 3 & 2 & 0\\
        0 & 1 & 2 & -1 & 7\\
        0 & 0 & 0 & 0 & 16\\
        0 & 0 & 0 & 0 & 0
    \end{pNiceMatrix}\]
    所以$\rank{\mat{A}}=3$. \sokka{C}
\end{solution}

\begin{choice}{A}
    设$n$阶矩阵$\mat{A}$非奇异($n\geq 2$),则
\begin{tasks}(4)
    \task $\ab(\mat{A^*})^*=\abs{\mat{A}}^{n-2}\mat{A}$
    \task $\ab(\mat{A^*})^*=\abs{\mat{A}}^{n-1}\mat{A}$
    \task $\ab(\mat{A^*})^*=\abs{\mat{A}}^{n+1}\mat{A}$
    \task $\ab(\mat{A^*})^*=\abs{\mat{A}}^{n+2}\mat{A}$
\end{tasks}
\end{choice}
\begin{solution}*
    由$\mat{A}^{-1}=\frac{\mat{A}^*}{\abs{\mat{A}}}$得$\mat{A}^*=\mat{A}^{-1}\abs{\mat{A}},\ \ab(\mat{A}^*)^{-1}=\frac{\mat{A}}{\abs{\mat{A}}},\ \abs{\mat{A}^*}^*=\abs{\mat{A}^*}\cdot\frac{\mat{A}}{\abs{\mat{A}}}=\abs{\mat{A}}^{n-2}\mat{A}$. \sokka{A}
\end{solution}

\newpage
\section{计算题（每题6分，共48分）}
\setcounter{prob}{0}
\begin{problem}
    求四阶行列式$\mathrm{D}=\begin{vmatrix}
        1 & -1 & 1 & x-1\\
        1 & -1 & x+1 & -1\\
        1 & x-1 & 1 & -1\\
        x+1 & -1 & 1 & -1
    \end{vmatrix}$.
\end{problem}
\begin{solution}
    \[
        \mathrm{D}=\begin{vmatrix}
            1 & -1 & 1 & x-1\\
            1 & -1 & x+1 & -1\\
            1 & x-1 & 1 & -1\\
            x+1 & -1 & 1 & -1
        \end{vmatrix}\\
        =x\begin{vmatrix}
            1 & -1 & 1 & x-1\\
            1 & -1 & x+1 & -1\\
            1 & x-1 & 1 & -1\\
            1 & -1 & 1 & -1
        \end{vmatrix}\\
        =x
        \begin{vmatrix}
            1 & 0 & 0 & x\\
            1 & 0 & x & 0\\
            1 & x & 0 & 0\\
            1 & 0 & 0 & 0
        \end{vmatrix}=x\cdot(-1)^{4+1}\cdot(-x^3)=x^4
    \]
\end{solution}

\begin{problem}
    求矩阵$\mat{A}=\begin{pmatrix}
        1 & 2 & 0 & 0\\
        3 & 4 & 0 & 0\\
        0 & 0 & 5 & 6\\
        0 & 0 & 7 & 8
    \end{pmatrix}$的逆矩阵.
\end{problem}
\begin{solution}
\columnseprule=.5pt
\begin{multicols*}{2}
    \begin{itemize}
        \item $\mat{A}_1=\begin{pmatrix}
            1 & 2\\
            3 & 4
        \end{pmatrix},\ \mat{A}_2=\begin{pmatrix}
            5 & 6\\
            7 & 8
        \end{pmatrix}$，则$\mat{A}=\begin{pmatrix}
            \mat{A}_1 & \mat{O}\\
            \mat{O} & \mat{A}_2
        \end{pmatrix}$.
        \item $\mat{A}_1^{-1}=-\frac12\begin{pmatrix}
            4 & -2\\
            -3 & 1
        \end{pmatrix},\ \mat{A}_2^{-1}=-\frac12\begin{pmatrix}
            8 & -6\\
            -7 & 5
        \end{pmatrix}$.
        \item 由逆矩阵的性质可知
        \[\mat{A}^{-1}=\begin{pmatrix}
            \mat{A}_1^{-1} & \mat{O}\\
            \mat{O} & \mat{A}_2^{-1}
        \end{pmatrix}=\begin{pNiceArray}{cc:cc}
            -2 & 1 & 0 & 0\\
            3/2 & -1/2 & 0 & 0\\
            \hdottedline
            0 & 0 & -4 & 3\\
            0 & 0 & 7/2 & -5/2
        \end{pNiceArray}\]
    \end{itemize}
\end{multicols*}
\end{solution}

\begin{problem}
    若矩阵$\mat{A}=\begin{pmatrix}
        3 & 1 & 1 & 4\\
        \lambda & 4 & 10 & 1\\
        1 & 7 & 17 & 3\\
        2 & 2 & 4 & 3
    \end{pmatrix}$
    经过初等行变换所得的行阶梯形矩阵有两个非零行，求$\lambda$的值.
\end{problem}
\begin{solution}
    \vspace{-1em}
    由题意得$\rank{\mat{A}}=2$，所以$\abs{\mat{A}}$中任意三阶子式均为$0$，即$\begin{vmatrix}
        3 & 1 & 1\\
        \lambda & 4 & 10\\
        1 & 7 & 17
    \end{vmatrix}=0,\ \lambda=0$.
\end{solution}

\begin{problem}
    设$\mathrm{D}=\begin{vmatrix}
        2 & 1 & 3 & 4\\
        1 & 0 & 2 & 3\\
        1 & 5 & 2 & 1\\
        -1 & 1 & 5 & 2
    \end{vmatrix}$
    ，求$\mathrm{A}_{13}+\mathrm{A}_{23}+\mathrm{A}_{43}$.
\end{problem}
\begin{solution}
\begin{itemize}
    \item 记$\mathrm{D}_1=\begin{vNiceMatrix}
        \CodeBefore
            \tikz \fill [H6,opacity=.2] (1-|3) |- (5-|4) |- (1-|4) |- cycle;
        \Body
        2 & 1 & 1 & 4\\
        1 & 0 & 1 & 3\\
        1 & 5 & 0 & 1\\
        -1 & 1 & 1 & 2
    \end{vNiceMatrix}$，则$\mathrm{A}_{13}+\mathrm{A}_{23}+\mathrm{A}_{43}$的值等价于$\mathrm{D}_1$按第三列展开
    \item $\mathrm{A}_{13}+\mathrm{A}_{23}+\mathrm{A}_{43}=\mathrm{D}_1\xlongequal[r_1-r_4]{r_3-5r_4}\begin{vmatrix}
        3 & 0 & 0 & 2\\
        1 & 0 & 1 & 3\\
        6 & 0 & -5&-9\\
        -1& 1 & 1 & 2
    \end{vmatrix}=(-1)^{4+2}\begin{vmatrix}
        3 & 0 & 2\\
        1 & 1 & 3\\
        6 &-5 & 9
    \end{vmatrix}=-4$.
\end{itemize}
\end{solution}

\begin{problem}
    已知线性变换$\begin{cases}
        y_1=x_1+2x_2-x_3\\
        y_2=3x_1+4x_2-2x_3\\
        y_3=5x_1-4x_2+x_3
    \end{cases}$
    ，求从变量$y_1,\ y_2,\ y_3$到变量$x_1,\ x_2,\ x_3$的线性变换.
\end{problem}
\begin{solution}
    记$\mat{A}=\begin{pNiceArray}{ccc:ccc}
        1 & 2 & -1 & 1 & 0 & 0\\
        3 & 4 & -2 & 0 & 1 & 0\\
        5 & -4 & 1 & 0 & 0 & 1
    \end{pNiceArray},\ \mat{A}\overset{r}{\sim}
    \begin{pNiceArray}{ccc:ccc}
        1 & 0 & 0 & -2 & 1 & 0\\
        0 & 1 & 0 & -13/2 & 3 & -1/2\\
        0 & 0 & 1 & -16 & 7 & -1
    \end{pNiceArray}$，所以得对应的线性变换为
    \[\mat{X}=\mat{A}^{-1}\mat{Y},\ \begin{cases}
        x_1=-2y_1+y_2\\
        x_2=-\frac{13}{2}y_1+3y_2-\frac12y_3\\
        x_3=-16y_1+7y_2-y_3
    \end{cases}\]
\end{solution}

\begin{problem}
    设$\mat{A}$是$n$阶可逆矩阵，将$A$中第$j$行所有元的$k$倍加到第$i$行对应的元素上去，得到的矩阵记为$\mat{B}$，求$\mat{AB}^{-1}$.
\end{problem}
\begin{solution}*
    由题意得$\rank{\mat{A}}=n$，由于$\mat{A}\overset{r}{\sim}\mat{B}$，所以$\rank{\mat{B}}=n$，$\mat{B}$可逆. 
    \[\mat{B}=\mat{E}\ab(i,j(k))\mat{A},\ \mat{B}^{-1}=(\mat{E}(i,j(k))\cdot\mat{A})^{-1}=\mat{A}^{-1}\cdot\mat{E}^{-1}(i,j(k))=\mat{A}^{-1}\mat{E}(i,j(-k))\]
    所以$\mat{AB}^{-1}=\mat{AA}^{-1}\mat{E}(i,j(-k))=\mat{E}(i,j(-k))$.
\end{solution}

\begin{problem}
    若齐次线性方程组$\begin{cases}
        x_1-2x_2+3\lambda x_3=0\\
        x_1-2\lambda x_2+3x_3=0\\
        \lambda x_1-2x_2+3x_3=0
    \end{cases}$
    有非零解，则$\lambda$的取值为？
\end{problem}
\begin{solution}
    该方程组有非零解，记该方程为$\mat{A}\bm{x}=0$. 由于$\rank{\mat{A}}<3$，所以$\abs{\mat{A}}=0$.
    \[\abs{\mat{A}}=\begin{vmatrix}
        1 & -2 & 3\lambda\\
        1 & -2\lambda & 3\\
        \lambda & -2 & 3
    \end{vmatrix}
    =6(\lambda-1)^2(\lambda+2)=0\]
    得$\lambda$的取值为$\lambda=1$或$\lambda=-2$.
\end{solution}

\begin{problem}
已知$\mat{A}=\begin{pmatrix}
    3 & 1 & 1 & 1\\
    1 & 4 & 3 & 4\\
    1 & 4 & 11 & 16\\
    1 & 8 & 27 & 66
\end{pmatrix}$
，且$\mat{AX}+4\mat{E}=\mat{A}^2+2\mat{X}$，求矩阵$\mat{X}$.
\end{problem}
\begin{solution}
    由条件$(\mat{A}-2\mat{E})\mat{X}=\mat{A}^2-4\mat{E}=(\mat{A}-2\mat{E})(\mat{A}+2\mat{E})$可得
    \[\abs{\mat{A}-2\mat{E}}=\abs{(\mat{A}-2\mat{E})^{-1}}=\begin{vmatrix}
        1 & 1 & 1 & 1\\
        1 & 2 & 4 & 8\\
        1 & 3 & 9 & 27\\
        1 & 4 & 16 & 64
    \end{vmatrix}=\prod_{4\geq i>j\geq 1}(x_i-x_4)=12\neq 0\]
    所以$\abs{\mat{A}-2\mat{E}}$可逆. 等式两边左乘$(\mat{A}-2\mat{E})^{-1}$得$\mat{X}=\mat{A}+2\mat{E}=\begin{pmatrix}
        5 & 1 & 1 & 1\\
        1 & 6 & 3 & 4\\
        1 & 4 & 13 & 16\\
        1 & 8 & 27 & 68
    \end{pmatrix}$.
\end{solution}

\newpage
\section{综合题（共10分）}
\setcounter{prob}{0}
\begin{problem}
设线性方程组$\begin{cases}
    x_1+x_2+2x_3+3x_4=1\\
    x_1+3x_2+6x_3+x_4=3\\
    3x_1-x_2-ax_3+15x_4=3\\
    x_1-5x_2-10x_3+12x_4=b
\end{cases}$，
问$a$与$b$各取何值时，线性方程组有唯一解？无解？无穷多解？并求出无穷多解时的通解.
\end{problem}
\begin{solution}
    增广矩阵$(\mat{A},\bm{b})=\begin{pNiceArray}{cccc:c}
        1 & 1 & 2 & 3 & 1\\
        1 & 3 & 6 & 1 & 3\\
        3 & -1 & -a & 15 & 3\\
        1 & -5 & -10 & 12 & b
    \end{pNiceArray}$，其行变换为
    \[(\mat{A},\bm{b})\overset{r}{\sim}\begin{pmatrix}
        1 & 1 & 2 & 3 & 1\\
        0 & 2 & 4 & -2& 2\\
        0 &-6 &-12& 9 &b-1\\
        0 &-4 &-a-b&6 & 0
    \end{pmatrix}\overset{r}{\sim}
    \begin{pmatrix}
        1 & 1 & 2 & 3 & 1\\
        0 & 1 & 2 &-1 & 1\\
        0 & 0 & 0 & 3&b+5\\
        0 & 0 &2-a& 2 & 4
    \end{pmatrix}\overset{r}{\sim}
    \begin{pmatrix}
        1 & 1 & 2 & 3 & 1\\
        0 & 1 & 2 &-1 & 1\\
        0 & 0 &2-a& 2 & 4\\
        0 & 0 & 0 & 3& b+5
    \end{pmatrix}\]
    \begin{itemize}
        \item 有唯一解：$\rank{\mat{A}}=\rank{\mat{A},\bm{b}}=n=4$，此时$a\neq 2$.
        \item 无解：$\rank{\mat{A}}<\rank{\mat{A},\bm{b}}$，此时$a=2,\ \rank{\mat{A}}=3,\ \rank{\mat{A},\bm{b}}=4$. 故$b+5\neq 6,\ b\neq 1$.
        \item 无穷多解：$\rank{\mat{A}}=\rank{\mat{A},\bm{b}}<4$，此时$a=2,\ b=1$.
    \end{itemize}
\end{solution}

\newpage
\section{综合题（共6分）}
\setcounter{prob}{0}
\begin{problem}
    设$n$阶矩阵$\mat{A}$和$\mat{B}$满足条件$\mat{BA}-3\mat{B}-2\mat{A}=0$.
\begin{enumerate}
    \item 证明$\mat{B}-2\mat{E}$为可逆矩阵.
    \item 已知$\mat{B}=\begin{pmatrix}
        2 & -3 & 0\\
        2 & 2 & 0\\
        0 & 0 & 3
    \end{pmatrix}$，求矩阵$\mat{A}$.
\end{enumerate}
\end{problem}
\begin{solution}
\begin{enumerate}
    \item \begin{proof}
        \[(\mat{B}-2\mat{E})\mat{A}=3\mat{B}=3(\mat{B}-2\mat{E})+6\mat{E}\Rightarrow(\mat{B}-2\mat{E})(\mat{A}-3\mat{E})=6\mat{E},\ \ab(\frac{\mat{B}-2\mat{E}}{6})(\mat{A}-3\mat{E})=6\mat{E}\]
        由逆矩阵可知$\mat{B}-2\mat{E}$可逆，且$(\mat{B}-2\mat{E})^{-1}=\frac{\mat{A}-3\mat{E}}{6}$.
    \end{proof}
    \item $\mat{B}-2\mat{E}=\begin{pmatrix}
        0 & -3 & 0\\
        2 & 0 & 0\\
        0 & 0 & 1
    \end{pmatrix}$
    \[
    \begin{pNiceArray}{ccc:ccc}
        0 & -3 & 0 & 1 & 0 & 0\\
        2 & 0 & 0 & 0 & 1 & 0\\
        0 & 0 & 1 & 0 & 0 & 1
    \end{pNiceArray}
    \xlongrightarrow[\substack{r_1\times\frac12\\r_2\times-\frac13}]{r_2\leftrightarrow r_1}
    \begin{pNiceArray}{ccc:ccc}
        1 & 0 & 0 & 0 & \frac12 & 0\\
        0 & 1 & 0 & -\frac13 & 0 & 0\\
        0 & 0 & 1 & 0 & 0 & 1
    \end{pNiceArray}
    \]
    所以$(\mat{B}-2\mat{E})^{-1}=\frac{\mat{A}-3\mat{E}}{6}=\begin{pmatrix}
        0 & \frac12 & 0\\
        -\frac13 & 0 & 0\\
        0 & 0 & 1
    \end{pmatrix},\ \mat{A}=\begin{pmatrix}
        3 & 3 & 0\\
        -2 & 3 & 0\\
        0 & 0 & 9
    \end{pmatrix}$.
\end{enumerate}
\end{solution}


\vfill
\rightline{/*Powered by \pkg{LiteSolution} \textsf\LaTeX~Template. Designed by WeYoung \href{https://ctan.org/home/xiamyphys}{Axia}.*/}